\section{Benchmarks}
\label{sec:benchmarks}

We benchmark the C and Jasmin implementations on three Intel microarchitectures: Skylake, Coffee Lake, and Raptor Lake, with Hyper-Threading and Turbo 
Boost disabled. For the hybrid Raptor Lake processor, we pin the benchmark to a performance core. All code is compiled using \texttt{GCC 12.2.0}.

We compare the \textsc{Kaiburr} variants with FrodoKEM and HQC. We use the SHAKE-based variant of FrodoKEM as the
primary comparison baseline. Both FrodoKEM-SHAKE and \textsc{Kaiburr} use SHAKE-based symmetric primitives.
However, FrodoKEM-AES can take advantage of dedicated AES instructions on our Intel processors, which introduces an 
implementation-specific hardware advantage unrelated to the lattice-based KEM operations we are looking to compare.\footnote{Additionally,
we do not report a direct cycle-count comparison with KyFrog~\cite{kyfrog}, the KEM closest to \textsc{Kaiburr} in principle, since no
public hand-optimized implementation of the scheme exists.}
We report only FrodoKEM-640, the smallest standardized FrodoKEM parameter set, since it is already substantially 
more expensive than \textsc{Kaiburr}. The larger FrodoKEM parameter sets have even higher computational costs and 
do not provide a more informative performance baseline.

All three KEMs are evaluated using a common benchmark driver adapted from the implementation used in~\cite{fvk}. 
We first perform $1000$ untimed warm-up iterations, followed by $10001$ timed iterations. We report the median 
execution time and interquartile range ($Q_1$--$Q_3$). This procedure is repeated three times, and we report the 
results from the run whose median is itself the median of the three runs, together with its corresponding quartiles. 
Each iteration verifies that the decapsulated shared secret matches the encapsulated shared secret.

Randomness comes from Linux \texttt{getrandom}. FrodoKEM and \textsc{Kaiburr}
call it within \texttt{keypair} and \texttt{enc}, so that syscall is inside
the measured region. HQC instead seeds a SHAKE-256 PRNG once, which we seed
from \texttt{getrandom} before timing begins.

The benchmark results are reported separately for each processor. Results for the Intel Core i7-6700K,
Intel Core i7-8700K, and Intel Core i9-13900K are given in Tables~\ref{tab:6700k}, \ref{tab:8700k}, and
\ref{tab:13900k}, respectively.


\begin{table*}[t]
\centering
\small
\caption{Benchmark results in cycles on an Intel Core i7-6700K. We report only the AVX2 implementations.
The reference implementations are, on average, 545.7 percent slower for \textsc{Kaiburr} C and 637.8 percent
 slower for \textsc{Kaiburr} Jasmin, measured using the geometric mean. Values are reported as median (Q1--Q3).}
\label{tab:6700k}
\begin{tabular}{llrrr}
\toprule
KEM & Impl. & KeyGen & Enc & Dec \\
\midrule
kaiburr-1792 & C/asm &
\makecell[r]{156186\\(155490--157822)} &
\makecell[r]{160062\\(159576--161914)} &
\makecell[r]{171888\\(171050--173356)} \\
kaiburr-1792 & Jasmin &
\makecell[r]{146528\\(145948--147226)} &
\makecell[r]{148044\\(147606--148788)} &
\makecell[r]{151990\\(151322--152728)} \\
\addlinespace
kaiburr-4608 & C/asm &
\makecell[r]{809842\\(807208--813058)} &
\makecell[r]{817844\\(815224--821062)} &
\makecell[r]{845458\\(842838--848794)} \\
kaiburr-4608 & Jasmin &
\makecell[r]{780574\\(779084--782266)} &
\makecell[r]{781302\\(779932--782906)} &
\makecell[r]{796676\\(795256--798396)} \\
\addlinespace
kaiburr-6144 & C/asm &
\makecell[r]{1361664\\(1358070--1366138)} &
\makecell[r]{1365370\\(1361692--1369542)} &
\makecell[r]{1407632\\(1403884--1411932)} \\
kaiburr-6144 & Jasmin &
\makecell[r]{1308772\\(1306774--1311182)} &
\makecell[r]{1309238\\(1307196--1312060)} &
\makecell[r]{1324338\\(1322338--1326730)} \\
\midrule
FrodoKEM-640 & \cite{frodoKEM-github} &
\makecell[r]{3860180\\(3859268--3868164)} &
\makecell[r]{4077768\\(4076966--4087064)} &
\makecell[r]{4044176\\(4043428--4053764)} \\
\midrule
HQC-1 & \cite{hqc-github} &
\makecell[r]{88162\\(88006--89002)} &
\makecell[r]{172786\\(172618--173672)} &
\makecell[r]{405674\\(405402--406028)} \\
HQC-3 & \cite{hqc-github} &
\makecell[r]{209986\\(209758--210708)} &
\makecell[r]{407810\\(407430--408654)} &
\makecell[r]{858366\\(857904--859064)} \\
HQC-5 & \cite{hqc-github} &
\makecell[r]{433948\\(433570--434640)} &
\makecell[r]{847696\\(846974--848576)} &
\makecell[r]{1712396\\(1711782--1713680)} \\
\bottomrule
\end{tabular}
\end{table*}


\begin{table*}[t]
\centering
\small
\caption{Benchmark results in cycles on an Intel Core i7-8700K. We report only the AVX2 implementations. 
The reference implementations are, on average, 540.8 percent slower for \textsc{Kaiburr} C and 653.8 percent
 slower for \textsc{Kaiburr} Jasmin, measured using the geometric mean. Values are reported as median (Q1--Q3).}
\label{tab:8700k}
\begin{tabular}{llrrr}
\toprule
KEM & Impl. & KeyGen & Enc & Dec \\
\midrule
kaiburr-1792 & C/asm &
\makecell[r]{156156\\(155579--157778)} &
\makecell[r]{159719\\(159235--161416)} &
\makecell[r]{173015\\(172106--174398)} \\
kaiburr-1792 & Jasmin &
\makecell[r]{143857\\(143402--144559)} &
\makecell[r]{144987\\(144656--145831)} &
\makecell[r]{150435\\(149704--151098)} \\
\addlinespace
kaiburr-4608 & C/asm &
\makecell[r]{814685\\(812084--817870)} &
\makecell[r]{821346\\(818719--824533)} &
\makecell[r]{848608\\(845755--852000)} \\
kaiburr-4608 & Jasmin &
\makecell[r]{762767\\(761322--764429)} &
\makecell[r]{765475\\(764033--767156)} &
\makecell[r]{781079\\(779648--782776)} \\
\addlinespace
kaiburr-6144 & C/asm &
\makecell[r]{1367405\\(1363808--1371615)} &
\makecell[r]{1376530\\(1372953--1380951)} &
\makecell[r]{1411130\\(1407567--1415488)} \\
kaiburr-6144 & Jasmin &
\makecell[r]{1293036\\(1291067--1295424)} &
\makecell[r]{1295674\\(1293713--1298070)} &
\makecell[r]{1313175\\(1311174--1315640)} \\
\midrule
FrodoKEM-640 & \cite{frodoKEM-github} &
\makecell[r]{3851932\\(3851203--3859217)} &
\makecell[r]{4067952\\(4067235--4075634)} &
\makecell[r]{4044333\\(4043732--4053215)} \\
\midrule
HQC-1 & \cite{hqc-github} &
\makecell[r]{90229\\(90070--90974)} &
\makecell[r]{172063\\(171888--172892)} &
\makecell[r]{404274\\(404018--404596)} \\
HQC-3 & \cite{hqc-github} &
\makecell[r]{215176\\(214956--215857)} &
\makecell[r]{404687\\(404324--405517)} &
\makecell[r]{857576\\(857048--858323)} \\
HQC-5 & \cite{hqc-github} &
\makecell[r]{439786\\(439458--440546)} &
\makecell[r]{838353\\(837433--839669)} &
\makecell[r]{1698485\\(1697851--1699766)} \\
\bottomrule
\end{tabular}
\end{table*}

\begin{table*}[t]
\centering
\small
\caption{Benchmark results in cycles on an Intel Core i9-13900K. We report only the AVX2 implementations. The reference implementations are, 
on average, 498.2 percent slower for \textsc{Kaiburr} C and 557.8 percent slower for \textsc{Kaiburr} Jasmin, measured using the geometric mean.
Values are reported as median (Q1--Q3).}
\label{tab:13900k}
\begin{tabular}{llrrr}
\toprule
KEM & Impl. & KeyGen & Enc & Dec \\
\midrule
kaiburr-1792 & C/asm &
\makecell[r]{121724\\(120872--122872)} &
\makecell[r]{125204\\(124310--126302)} &
\makecell[r]{134816\\(134058--136082)} \\
kaiburr-1792 & Jasmin &
\makecell[r]{141364\\(140882--142142)} &
\makecell[r]{143058\\(142620--143842)} &
\makecell[r]{148502\\(147944--149212)} \\
\addlinespace
kaiburr-4608 & C/asm &
\makecell[r]{648064\\(645678--650674)} &
\makecell[r]{654366\\(651728--657234)} &
\makecell[r]{676840\\(674446--679554)} \\
kaiburr-4608 & Jasmin &
\makecell[r]{747376\\(745792--749120)} &
\makecell[r]{749472\\(747868--751250)} &
\makecell[r]{764490\\(762874--766280)} \\
\addlinespace
kaiburr-6144 & C/asm &
\makecell[r]{1095102\\(1091646--1098922)} &
\makecell[r]{1097334\\(1093566--1101348)} &
\makecell[r]{1127588\\(1124030--1131712)} \\
kaiburr-6144 & Jasmin &
\makecell[r]{1272644\\(1270318--1275208)} &
\makecell[r]{1272480\\(1270246--1274906)} &
\makecell[r]{1285796\\(1283420--1288436)} \\
\midrule
FrodoKEM-640 & \cite{frodoKEM-github} &
\makecell[r]{3112706\\(3107260--3124744)} &
\makecell[r]{3239826\\(3233936--3249204)} &
\makecell[r]{3195408\\(3187982--3208002)} \\
\midrule
HQC-1 & \cite{hqc-github} &
\makecell[r]{72578\\(71902--73220)} &
\makecell[r]{139896\\(138680--141026)} &
\makecell[r]{312692\\(310798--314622)} \\
HQC-3 & \cite{hqc-github} &
\makecell[r]{179494\\(178058--180668)} &
\makecell[r]{332806\\(329936--335788)} &
\makecell[r]{668334\\(663984--672664)} \\
HQC-5 & \cite{hqc-github} &
\makecell[r]{377670\\(376224--380158)} &
\makecell[r]{705038\\(700890--709262)} &
\makecell[r]{1357572\\(1350500--1364182)} \\
\bottomrule
\end{tabular}
\end{table*}

\textsc{Kaiburr} is faster than FrodoKEM-640 across every
\textbf{operation}, \textbf{parameter set}, \textbf{implementation}, and \textbf{machine}.
 kaiburr-1792 is 21.5--28.1x faster, kaiburr-4608 is 4.2--5.3x faster, and kaiburr-6144 is
2.5--3.1x faster than FrodoKEM-640.

Comparing the \textsc{Kaiburr} C/asm and Jasmin implementations, the faster implementation depends on
the microarchitecture. On the i7-6700K and i7-8700K, Jasmin is faster than C/asm for every parameter set
and every operation, by 3.6--11.6\% and 5.4--13.1\%, respectively. On the i9-13900K, this reverses and C/asm is faster than Jasmin for every
parameter set and every operation, by 9.2--14.0\%.

Across parameter sets, C/asm KeyGen cycle counts scale by a factor of 5.19 (6700K) / 5.22 (8700K) / 5.32 (13900K) from
kaiburr-1792 ($k=7$) to kaiburr-4608 ($k=18$), and by a factor of 1.68 (6700K) / 1.68 (8700K) / 1.69 (13900K) from
kaiburr-4608 to kaiburr-6144 ($k=24$). For reference, $k$ itself grows by factors of 2.57 and 1.33 over
these two steps, and $k^2$ by factors of 6.61 and 1.78.

kaiburr-1792 is faster than HQC-5 in every operation on all three machines: by 64.0--90.0\% on the
6700K (KeyGen, Enc, Dec), 64.5--89.8\% on the 8700K, and 67.8--90.1\% on the 13900K. Comparing each \textsc{Kaiburr} parameter set to the
HQC parameter set of the same index (kaiburr-1792/HQC-1, kaiburr-4608/HQC-3, kaiburr-6144/HQC-5), HQC's
KeyGen is faster than the corresponding \textsc{Kaiburr} KeyGen by a factor of 1.77, 3.86, and 3.14 on
the 6700K, 1.73, 3.79, and 3.11 on the 8700K, and 1.68, 3.61, and 2.90 on the 13900K.

Interquartile ranges are narrow relative to the medians across all measurements, typically within
1--2\% of the median (e.g., for kaiburr-1792 C/asm, 1.3--1.5\% on the 6700K and 1.3--1.4\% on the 8700K).

We additionally benchmark the ACC implementation described in Section~\ref{sec:implementation} on
Pavona's cycle-accurate RTL simulator for the accelerator. Because the simulator is deterministic
for a fixed random seed, repeated runs of the same test are identical in cycles. Therefore unlike the CPU benchmarks
above, we report a single exact cycle count per operation rather than a median and interquartile
range. We measure each KEM operation within a single keypair--encapsulate--decapsulate
roundtrip and record cumulative execution statistics after each operation. The individual costs
are then recovered by subtraction:
$\mathrm{KeyGen}=M_1$, $\mathrm{Enc}=M_2-M_1$, and
$\mathrm{Dec}=M_3-M_2$. Results are shown in Table~\ref{tab:acc}.

\begin{table}[t]
\centering
\caption{Benchmark results in cycles for the ACC implementation on Pavona's cycle-accurate RTL
simulator. Values are exact cycle counts.}
\label{tab:acc}
\begin{tabular}{lrrr}
\toprule
KEM & KeyGen & Enc & Dec \\
\midrule
kaiburr-1792 & 240816 & 239348 & 257474 \\
kaiburr-4608 & 1453490 & 1424342 & 1467109 \\
kaiburr-6144 & 2505793 & 2449864 & 2506070 \\
\bottomrule
\end{tabular}
\end{table}

Table~\ref{tab:acc-profile} shows the three functions accounting for the largest share of cycles
in the same roundtrip benchmark, obtained from the simulator's per-function execution statistics.
Across all three parameter sets, matrix generation (\texttt{poly\_gen\_matrix}) dominates,
consuming 73--84\% of total cycles.

\begin{table}[t]
\centering
\caption{The three functions with the largest share of total cycles in the ACC roundtrip
benchmark, ranked 1st to 3rd, obtained from the simulator's per-function execution statistics.
All values are percentages of each variant's total roundtrip cycles. Cell shading is scaled to
the largest value in the table.}
\label{tab:acc-profile}
\begin{tabular}{lrrr}
\toprule
KEM & 1st & 2nd & 3rd \\
\midrule
kaiburr-1792 &
\heat{72.60}{83.56}{\texttt{poly\_gen\_matrix} 72.60} &
\heat{8.08}{83.56}{\texttt{basemul\_acc} 8.08} &
\heat{6.26}{83.56}{\texttt{f4} 6.26} \\
kaiburr-4608 &
\heat{80.78}{83.56}{\texttt{poly\_gen\_matrix} 80.78} &
\heat{9.23}{83.56}{\texttt{basemul\_acc} 9.23} &
\heat{3.96}{83.56}{\texttt{f6} 3.96} \\
kaiburr-6144 &
\heat{83.56}{83.56}{\texttt{poly\_gen\_matrix} 83.56} &
\heat{9.57}{83.56}{\texttt{basemul\_acc} 9.57} &
\heat{2.05}{83.56}{\texttt{f8} 2.05} \\
\bottomrule
\end{tabular}
\end{table}

\texttt{poly\_gen\_matrix} expands the public seed into the uniformly random matrix
$\mathbf{A}$ via rejection sampling on SHAKE-128 output. Since it dominates the ACC cycle counts,
we break down its costs further using the simulator's per-instruction execution statistics.
For kaiburr-6144, this instruction-level breakdown of \texttt{poly\_gen\_matrix} shows that
rejection sampling is responsible for the vast majority of its
cost. Only 464{,}022 of its 6{,}235{,}212 cycles (7.4\%, or 6.2\% of the entire benchmark) are
spent stalled on \texttt{bn.wsrr} waiting for the hardware SHAKE-128 core to produce the next
squeeze block. The remaining 92.6\% is software cost of the rejection-sampling loop itself. 83.6\%
of \texttt{poly\_gen\_matrix}'s cycles (69.9\% of the whole benchmark) are retired arithmetic
instructions (\texttt{bn.and}, \texttt{bn.cmp}, and \texttt{bn.sel}) performing the
mask/compare/accumulate operations of rejection sampling, and a further 8.5\% (7.1\% of the whole
benchmark) comes from branch overhead (\texttt{bne} and \texttt{beq}) in the accept/reject loop.
We repeated this instruction-level breakdown for kaiburr-1792 and kaiburr-4608 and found the same
pattern. The hardware-Keccak-wait share stays flat (5.5\% and 6.0\% of each
benchmark's total cycles, respectively, versus 6.2\% for kaiburr-6144), while the rejection-sampling
arithmetic share grows with the parameter set (60.6\% and 67.6\%, versus 69.9\% for kaiburr-6144).
Table~\ref{tab:acc-keccak} gives the full comparison.

\begin{table}[t]
\centering
\caption{Cycle breakdown of \texttt{poly\_gen\_matrix} across parameter sets, obtained from the
simulator's per-instruction execution statistics. All values are percentages of each variant's
total roundtrip cycles. Cell shading is scaled to the largest value in the table.}
\label{tab:acc-keccak}
\begin{tabular}{lrrrr}
\toprule
KEM & \makecell[r]{Rejection\\sampling} & \makecell[r]{Branch\\overhead}
    & \makecell[r]{Keccak\\wait} & \makecell[r]{Memory\\latency} \\
\midrule
kaiburr-1792 &
\heat{60.6}{69.9}{60.6} & \heat{6.2}{69.9}{6.2} & \heat{5.5}{69.9}{5.5} & \heat{0.3}{69.9}{0.3} \\
kaiburr-4608 &
\heat{67.6}{69.9}{67.6} & \heat{6.8}{69.9}{6.8} & \heat{6.0}{69.9}{6.0} & \heat{0.4}{69.9}{0.4} \\
kaiburr-6144 &
\heat{69.9}{69.9}{69.9} & \heat{7.1}{69.9}{7.1} & \heat{6.2}{69.9}{6.2} & \heat{0.4}{69.9}{0.4} \\
\bottomrule
\end{tabular}
\end{table}
